\rtmtight
\begin{tblr}{width=\textwidth,colspec={|Q[c,m,2.1cm]|X[l,m]|},hlines,vlines,hspan=minimal,rowsep=1.5pt,colsep=3pt,row{1}={bg=headgray,font=\bfseries}}
\SetCell{c,m}\hfil\logo\hfil & \textbf{POLITEKNIK NEGERI MALANG}\par\textbf{JURUSAN TEKNOLOGI INFORMASI}\par\textbf{PROGRAM STUDI: D4 TEKNIK INFORMATIKA} \\
\SetCell[c=2]{l} \textbf{RENCANA TUGAS MAHASISWA (RTM) DAN RUBRIK PENILAIAN} & \\
Nama Mata Kuliah & Sistem Informasi Manajemen \\
Kode Mata Kuliah & RTI253002 \quad \textbf{sks / Jam perminggu:} 2 SKS/4 jam \quad \textbf{Semester:} 3 \\
Dosen Pengampu & Tim Pengajar Program Studi D4 TI \\
\end{tblr}

\assessmentblock{Tugas Konsep dan Analisis SIM}{Case Method}{SCPMK0211-02001, SCPMK0211-02002}{Mahasiswa menyelesaikan tugas analisis konsep dasar sistem informasi manajemen sesuai Sub-CPMK dan konteks mata kuliah.}{Menganalisis kebutuhan organisasi, mengerjakan analisis terstruktur, mendokumentasikan evidence, dan mempresentasikan temuan.}{Laporan analisis, diagram sistem, dokumentasi proses bisnis, dan/atau bahan presentasi.}{Analisis konsep dan komponen sistem & 8 & Konsep sistem, informasi, dan komponen SIM diidentifikasi secara lengkap dan dikaitkan dengan kebutuhan organisasi nyata. & Sebagian besar komponen teridentifikasi tetapi hubungan dengan kebutuhan organisasi belum kuat. & Analisis tidak menyeluruh atau tidak terkait konteks organisasi. \\ \\
Pemodelan proses bisnis dan alur keputusan & 7 & Proses bisnis dan alur pengambilan keputusan dimodelkan secara sistematis dan konsisten dengan struktur SIM. & Model bisnis ada tetapi alur keputusan atau konsistensi dengan SIM belum optimal. & Pemodelan tidak mencerminkan proses bisnis atau alur keputusan yang logis. \\ \\
Kualitas dokumentasi dan presentasi & 5 & Laporan terstruktur, lengkap, terminologi benar, dan siap dipresentasikan. & Dokumentasi cukup namun beberapa bagian kurang rinci atau terminologi belum tepat. & Dokumentasi tidak lengkap atau tidak dapat mendukung presentasi. \\}

\assessmentblock{Kuis Konsep SIM (Kuis 1 dan Kuis 2)}{Kuis}{SCPMK0211-02001, SCPMK0607-02003}{Kuis tertulis untuk menguji pemahaman konsep dasar SIM dan penerapan SI strategis.}{Menjawab soal pilihan ganda dan/atau esai pendek terkait konsep SIM dan SI strategis.}{Jawaban tertulis kuis.}{Penguasaan konsep dasar SIM & 5 & Jawaban menunjukkan pemahaman mendalam tentang sistem, informasi, dan struktur SIM serta penerapannya. & Sebagian konsep dipahami namun penerapan pada kasus belum lengkap. & Jawaban hanya hafalan istilah tanpa penerapan konsep. \\ \\
Penerapan konsep SI strategis & 5 & SI strategis dan antar organisasi diidentifikasi dan dianalisis dengan tepat pada konteks kasus. & Identifikasi SI strategis cukup tepat namun analisis dampaknya belum menyeluruh. & Konsep SI strategis tidak dipahami atau tidak diterapkan dengan benar. \\}

\assessmentblock{CM Studi Kasus SI Strategis}{Case Method}{SCPMK0607-02003, SCPMK0802-02004}{Mahasiswa menganalisis studi kasus penerapan SI strategis pada organisasi nyata dan menyusun rekomendasi solusi.}{Menganalisis kasus SI organisasi, mengidentifikasi SI strategis yang diterapkan, merancang solusi integratif, dan mempresentasikan rekomendasi.}{Laporan studi kasus, matriks SI organisasi, rekomendasi solusi, dan bahan presentasi.}{Identifikasi dan analisis SI strategis & 12 & SI strategis diidentifikasi dengan tepat, dampaknya pada keunggulan kompetitif organisasi dianalisis secara mendalam dengan bukti kasus. & SI strategis teridentifikasi tetapi analisis dampak atau keterkaitannya dengan strategi organisasi belum kuat. & Identifikasi SI strategis tidak tepat atau tidak disertai analisis kontekstual. \\ \\
Perancangan solusi SIM terintegrasi & 10 & Solusi SIM dirancang secara komprehensif, mempertimbangkan integrasi antar sistem dan kebutuhan semua level organisasi. & Solusi dirancang namun integrasi antar komponen atau cakupan level organisasi masih terbatas. & Solusi tidak mempertimbangkan integrasi atau tidak sesuai kebutuhan organisasi. \\ \\
Kualitas presentasi dan argumentasi & 8 & Presentasi terstruktur, argumen didukung data kasus, dan mampu menjawab pertanyaan dengan tepat. & Presentasi cukup terstruktur namun argumentasi atau kemampuan menjawab pertanyaan masih terbatas. & Presentasi tidak terstruktur atau tidak dapat mempertahankan argumen. \\}

\assessmentblock{UTS Evaluasi Tengah Semester}{UTS berbasis analisis kasus}{SCPMK0211-02001, SCPMK0211-02002}{Evaluasi tengah semester untuk mengukur kemampuan analisis konsep SIM dan pemodelan sistem.}{Menganalisis kasus organisasi, merancang model SIM, dan mendokumentasikan solusi secara tertulis.}{Jawaban ujian tertulis dengan analisis dan model SIM.}{Analisis kebutuhan organisasi & 8 & Kebutuhan informasi dan proses bisnis organisasi diidentifikasi secara lengkap dan sistematis. & Sebagian kebutuhan teridentifikasi namun sistemasi analisis belum optimal. & Analisis kebutuhan tidak lengkap atau tidak relevan dengan konteks organisasi. \\ \\
Perancangan model konseptual SIM & 7 & Model konseptual SIM dirancang dengan notasi yang benar, mencakup semua komponen utama dan relasi antar komponen. & Model konseptual ada namun beberapa komponen atau relasi masih kurang tepat. & Model tidak mencerminkan pemahaman tentang struktur SIM. \\}

\assessmentblock{UAS Pengembangan SIM}{UAS berbasis proyek mini}{SCPMK0607-02003, SCPMK0802-02004}{Mahasiswa menyusun rencana pengembangan SIM menggunakan pendekatan SDLC untuk konteks organisasi yang diberikan.}{Menganalisis kebutuhan, memilih pendekatan SDLC yang tepat, menyusun rencana pengembangan lengkap, dan mendokumentasikan deliverable.}{Dokumen rencana pengembangan SIM, diagram alur SDLC, dan spesifikasi teknis awal.}{Ketepatan pendekatan SDLC & 7 & Pendekatan SDLC dipilih dengan justifikasi yang kuat berdasarkan karakteristik proyek dan kebutuhan organisasi. & Pendekatan SDLC dipilih namun justifikasi belum sepenuhnya meyakinkan. & Pendekatan SDLC tidak sesuai atau tidak disertai justifikasi yang memadai. \\ \\
Kelengkapan rencana pengembangan & 8 & Rencana pengembangan mencakup semua fase SDLC dengan deliverable, timeline, dan sumber daya yang jelas. & Rencana cukup lengkap namun beberapa fase atau deliverable masih kurang rinci. & Rencana tidak lengkap atau tidak dapat diimplementasikan. \\ \\
Kualitas spesifikasi teknis & 7 & Spesifikasi teknis awal ditulis dengan notasi yang benar, mencakup kebutuhan fungsional dan non-fungsional. & Spesifikasi teknis ada namun cakupan atau notasinya belum lengkap. & Spesifikasi teknis tidak ada atau tidak mencerminkan pemahaman rekayasa perangkat lunak. \\}

\begin{tblr}{width=\textwidth,colspec={|Q[l,m,3.2cm]|X[l,m]|},hlines,vlines,hspan=minimal,row{1-3}={bg=headgray},rowsep=1.5pt,colsep=3pt}
Jadwal Pelaksanaan: & Minggu 1--17 sesuai RPS. Setiap evaluasi memiliki RTM dan rubrik multi-kriteria. \\
Lain-Lain Yang Diperlukan: & LMS, repositori/proyek, perangkat lunak sesuai mata kuliah, dan perangkat presentasi. \\
Pustaka: & Mengacu pustaka utama dan pendukung pada bagian RPS. \\
\end{tblr}
